\begin{proof}
We first prove the one-step assertion.  Let \(h\) be classically
pluriharmonic on an ordinary neighbourhood \(\Omega\) of \(s(y)\), and put
\(\widetilde g:=s^*h\) on the ordinary neighbourhood
\(s^{-1}(\Omega)\) of \(y\).  By
\Cref{prop:p061-relative-positive-weighted-cone-descent}, after shrinking
there are a strict affinoid ordinary neighbourhood \(V\) of \(y\), a
weighted tube \(U_n(V)\subset\Omega\) whose ordinary interior contains
\(s(y)\), and the classically pluriharmonic function
\(g=\widetilde g|_V=s^*h|_V\) such that
\[
                         h|_{U_n(V)}=p^*g.
\tag{\(\diamondsuit\)}
\]
Property \({\rm PH}_{\mathbf R{\rm PL}}\) at \(y\), applied to
\(\widetilde g\) on the ordinary open set \(s^{-1}(\Omega)\), supplies an
ordinary open set \(y\in W\subset s^{-1}(\Omega)\) on which
\(\widetilde g\) is \(\mathbf R\)-PL.  The set
\[
       \operatorname{Int}_X(U_n(V))\cap p^{-1}(W)
\]
is an ordinary neighbourhood of \(s(y)\).  On this set, pullback stability
of \(\mathbf R\)-PL functions \cite[Lemma~5.4(1)]{GJR26} and
\((\diamondsuit)\) show that \(h=p^*(\widetilde g|_W)\) is
\(\mathbf R\)-PL.  This proves the
one-step closure without assuming that the base is quasi-smooth.

Each \(B_j\) in the statement is again a strict affinoid \(K\)-algebra, and
the map \(p_j\) with its zero section \(s_j\) has exactly the form required
by \Cref{prop:p061-relative-positive-weighted-cone-descent}; coefficients
of the \(j\)-th equations are allowed to be arbitrary elements of
\(B_{j-1}\).  Starting with \(x_0\), repeated application of the one-step
assertion therefore gives property \({\rm PH}_{\mathbf R{\rm PL}}\) at
\(x_1,\ldots,x_\ell\).

After the final shrinking, a classically pluriharmonic \(h\) is
\(\mathbf R\)-PL, and both \(h\) and \(-h\) are classically psh.  They are
semipositive by \cite[Theorem~4.3]{GR25Classical}; this is precisely
Gubler--Jell--Rabinoff harmonicity
\cite[Definition~7.2]{GJR26}.  Finally, if \(X_0\) is quasi-smooth near
\(x_0\), then \Cref{thm:p061-quasismooth-ph-rpl} gives the required initial
property.
\end{proof}
